\section{Data}\label{sec:data}
Summarize the DP2 context relevant to calibrations: instrument, filters, observing period/epochs, detector layout, available exposure types (bias/dark/flat/science), and definitions of calibration epochs.
Optionally add science context.

\subsection{LSSTCam and CCD status for DP2}\label{sec:data:lsstcam}

Describe LSSTCam

Describe the detector layout

Describe the electronic readout board (REB)

Describe the sequencer

Describe the filters.

Analyzed with updated sequencer (s3v2).

\textcolor{red}{to do:} insert map as a percentage of used pixels: dead, bad and the dead ccd are in one color, color for CCDs worst than the median and green for best CCDs.
Effective area: vignetted masked + defects.



\subsection{Calibration System (CalSys)}\label{sec:data:calsys}

LSST uses four calibration systems to construct calibrations 

\subsubsection{Camera Calibration Optical Bench}\label{sec:data:calsys:ccob}

To create optical scenes for characterization, the Camera is placed in-line with an optical scene projector called the Camera Calibration Optical Bench (CCOB) \citep{2024SPIE13103E..0WU}. The CCOB has two configurations.

The CCOB-wide beam is designed to project a uniform monochromatic light source onto the focal plane. Light is produced by one of six LEDs, each corresponding to a central wavelength in the LSST filter bandpass, is controlled by a custom LED driver and fed into an integrating sphere. The light from the integrating sphere is projected onto the focal plane with a minimal ND filter (10\%) attached to a C-mount lens with an adjustable f/stop. The CCOB-wide beam was used in run 6 and run 7.

\textbf{I mention run 6 and run 7 here - it is not properly introduced. Is it within scope to mention the EO-testing run at the Rubin Observatory L3 CR?}

The CCOB-narrow beam is designed to project narrow beams of monochromatic light onto specific locations of the focal plane to study spatial throughput, ghosting, and reflections. Light is produced by a tunable laser source on a X/Y and $\theta / \phi$ motor drives. Total throughput can be measured by combining different positions and angles of the CCOB-narrow beam. The CCOB-narrow beam was used in run 6 only.

One difference between run 6 and run 7 was that the f/stop of the lens was changed from 2.6 to 1.6 (fully open). This was done to try to reduce the effect of the `weather' and the `CMB pattern' two effects that we found in Run 6 at SLAC, and were found to be due to our projection setup (\citet{2024arXiv241113386B}). While changing the f/stop  did reduce the weather pattern, it also caused a much steeper illumination roll-off across the focal plane. We evaluated the weather pattern and illumination roll-off relative to Run 6.

To reduce the effects of the `weather' and `CMB' but retain uniform illumination across the focal plane, we installed a Lambertian diffuser to the CCOB-wide beam for Run 7. The diffuser used is a $60^o$ diffusing angle unmounted sheet from Edmunds Optics. We found that the diffuser greatly reduced the `weather', eliminated the CMB pattern, more uniformly illuminated the focal plane, with a penalty of decreasing the overall illumination by roughly 35\%. The diffuser was installed for all run 7 record runs using the CCOB Wide Beam.

\subsubsection{The Auxiliary Telescope (AuxTel)}\label{sec:data:calsys:auxtel}

Overview of AuxTel

What instruments are installed @ AuxTel 

What can we do with AuxTel data products

\subsubsection{Collimated Beam Projector}\label{sec:data:calsys:cbp}

The Collimated Beam Projector (CBP) is a tunable monochromatic calibration source installed inside the Rubin Observatory dome. The CBP is designed as a Schmidt reflector telescope run in reverse, with a fiber-fed tunable laser that covers the LSST bandpass from 300nm - 1100nm. The CBP supports five interchangeable masks for spot projection onto the LSSTCam focal-plane. The CBP produces spatially uniform projections across all LSSTCam CCDs through projecting light from the CBP onto different positions of the Simonyi Survey Telescope primary mirror. 

The interchangeable masks of the CBP help to serve different calibration objectives, including telescope response, filter throughput, and CCD quantum efficiency (QE). \textbf{Not sure how much to expand on this here. Also, this would be a good place to cite Nathan's CBP contribution to SPIE - under PUB-207.}

\subsubsection{Flat Field Illuminator}\label{sec:data:calsys:flat_field_illuminator}

The 

\paragraph{Broadband LEDs}\label{sec:data:calsys:flat_field_illuminator:LEDs}



\paragraph{Tunable Laser}\label{sec:data:calsys:flat_field_illuminator:laser}



% - any other instrumentation used for input datasets.

\subsubsection{Limitations to Calibrations during LSST Y1}\label{sec:data:calsys:limitations}

Describe any system-level constraints that influenced calibration strategy.

Describe the LSST Camera and any technical limitations (bad detectors/amps, CCD power cycling and its affect on existing calibs).

\subsection{Data selection and corresponding calibration epoch}\label{sec:data:selection}
Describe the data selection process for the DP2 dataset and calibration epoch. See \ref{fig:calibepoch}.

\begin{figure*}[t]
    \centering
    \fbox{\parbox{0.75\textwidth}{\centering\vspace{1.5in}\small\textit{Figure pending}\vspace{1.5in}}}
    % \includegraphics[width=0.8\textwidth]{figures/calibration_epoch.pdf}
    \caption{}
    \label{fig:calibepoch}
\end{figure*}

